\emph{Setting.} Theorem \claimref{synthesis.rung-field-on-the-lift-locus} supplies the coordinates: $Q(c) = \alpha^{\mathsf T}\beta$, $\mathsf{N}(c) = \tfrac12(|\alpha|^{2} + |\beta|^{2})$, the Hermitian form $G(c, d) = \operatorname{Re}(c^{*}\cdot d) = \tfrac12\operatorname{Re}(\alpha^{\dagger}\alpha' + \beta^{\dagger}\beta')$ whose quadratic form is $\mathsf{N}$, the action $\alpha \mapsto X\alpha$, $\beta \mapsto \bar{X}\beta$ of the slot mixing $U(2)$ (so that $\beta$ transforms in the conjugate of the fundamental, and $\bar{\beta}$ in the fundamental), the hypercharge moment map $m = \tfrac12(|\alpha|^{2} - |\beta|^{2}) = G(c, Kc)$ with $K = -2I_{4}Y$, the lift family $c_{\lambda}$ ($\alpha = (H^{2}/\lambda, 0)$, $\beta = (-\lambda, 0)$), its clause (1) (the charge-preserving points of $L$ are the $U(2)$-images of the family, with $|\alpha||\beta| = H^{2}$ there), its clause (2) (the gauge quadratic form of the Hermitian kinetic term at $c_{\lambda}$ is $\tfrac14 \mathsf{N}(c_{\lambda})P$), its clause (3) ($\mathsf{N}(dc_{\lambda}/d\eta) = \mathsf{N}(c_{\lambda})$, the neutral physical line spanned by $dc_{\lambda}/d\lambda$, the charged plane its $G$-orthogonal complement in the physical tangent space) and its clause (5) ($\mathsf{N}^{2} = |Q|^{2} + m^{2}$ exactly on the charge-preserving $c$, $\mathsf{N}^{2} > |Q|^{2} + m^{2}$ elsewhere). Theorem \claimref{synthesis.rung-unit-is-frame-boost} gives the frame as the boost $\lambda$ and the premise \claimref{synthesis.rung-unit-second-free-number} registers $v/H$ as that frame. Masses are read with the conventions of a real field $h$ with Lagrangian $\tfrac12(\partial h)^{2} - \tfrac12 m_{h}^{2}h^{2}$ and a real vector field with mass term $\tfrac12 M^{2}W^{2}$: from the stored gauge form, $\tfrac14 \mathsf{N}(g^{2}W_{1}^{2} + g^{2}W_{2}^{2} + (gW_{3} + g'B)^{2})$, $M_{W}^{2} = \tfrac12 g^{2}\mathsf{N}$ and $M_{Z}^{2} = \tfrac12(g^{2} + g'^{2})\mathsf{N}$; for a direction $d$ tangent to $L$ at the vacuum, the kinetic term along the curve $c + \epsilon d + \dots$ is $\mathsf{N}(d)(\partial\epsilon)^{2}$, so the canonically normalised field is $h = \sqrt{2\mathsf{N}(d)}\,\epsilon$ and the mass squared is the second $\epsilon$-derivative of the potential along a curve in $L$ divided by $2\mathsf{N}(d)$. Only the ratio of the scalar to the gauge masses is used, and it does not depend on the overall normalisation of the kinetic term provided the D-term is built from the same Hermitian form, as it is here. The check works with Gaussian rationals, $\Lambda = 3$ ($H^{2} = 1$), $g = 2/3$, $g' = 1/2$, three vacuum frames $\lambda^{*} = 1/2, 3, 10$ and two added quadratics $W$.

\emph{(1).} Put $u = \alpha$ and $u' = \bar{\beta}$; both transform in the fundamental of $U(2)$. By the first fundamental theorem for the unitary group (from memory: the polynomial invariants of $U(n)$ on several copies of $\C^{n}$ and their conjugates are generated by the Hermitian pairings), the invariants are generated by $u^{\dagger}u = |\alpha|^{2}$, $u'^{\dagger}u' = |\beta|^{2}$ and $u^{\dagger}u' = \overline{\alpha^{\mathsf T}\beta} = \bar{Q}$ with its conjugate $Q$. The check verifies the invariance of the three under five rational unitary matrices (a rotation, two diagonal phases and their products) on three sample vectors. The moment maps: $-I_{4}Y$ acts as $+\tfrac12$ on $\alpha$ and $-\tfrac12$ on $\beta$ (the conjugate action), so $\mu_{Y} = \tfrac12\cdot\tfrac12(|\alpha|^{2} - |\beta|^{2}) = \tfrac12 m$; and $-I_{4}T_{j}$ acts as $\tfrac12\sigma_{j}$ on $\alpha$ and as $-\tfrac12\bar{\sigma}_{j}$ on $\beta$, so $\mu_{j} = \tfrac14(\alpha^{\dagger}\sigma_{j}\alpha - \beta^{\dagger}\bar{\sigma}_{j}\beta) = \tfrac14(u^{\dagger}\sigma_{j}u - u'^{\dagger}\sigma_{j}u')$, since $\beta^{\dagger}\bar{\sigma}_{j}\beta = \overline{\beta^{\mathsf T}\sigma_{j}\bar{\beta}} = u'^{\dagger}\sigma_{j}u'$ is real. With the Fierz identity $\sum_{j}(u^{\dagger}\sigma_{j}u)(u'^{\dagger}\sigma_{j}u') = 2|u^{\dagger}u'|^{2} - |u|^{2}|u'|^{2}$ and $\sum_{j}(u^{\dagger}\sigma_{j}u)^{2} = |u|^{4}$,
\[
\sum_{j}\mu_{j}^{2} = \tfrac{1}{16}\bigl((|u|^{2} + |u'|^{2})^{2} - 4|u^{\dagger}u'|^{2}\bigr) = \tfrac14(\mathsf{N}^{2} - |Q|^{2}),
\]
verified exactly on the samples. On $L$, $|Q| = H^{2}$, so the invariants reduce to $(\mathsf{N}, m)$; the region they fill is $\mathsf{N}^{2} - m^{2} \ge H^{4}$ by clause (5) of the stored theorem (equality on the charge-preserving points, strict inequality elsewhere), and $\mathsf{N}^{2} - m^{2} = |\alpha|^{2}|\beta|^{2}$ identically, so the boundary is $|\alpha||\beta| = H^{2}$, the lift family's orbits by clause (1). On the family $\alpha$ and $\bar{\beta}$ are parallel, so $\mu_{1} = \mu_{2} = 0$ and $\mu_{3} = \tfrac14(|\alpha|^{2} - |\beta|^{2}) = \tfrac12 m$ (checked at four frames); the unbroken generator $Y - T_{3}$ has vanishing moment map there, as it must. The check also exhibits a point of $L$ off the family with $\mathsf{N}^{2} - m^{2} > H^{4}$ and all four moment maps nonzero (a trivial stabiliser). Charge conservation: the stabiliser of a point of $L$ is a circle exactly on the boundary (clause (1)), so a minimum of an invariant potential conserves a $U(1)$ iff it lies on the boundary; a potential $f(\mathsf{N})$ alone has as minimum set the curve $\mathsf{N} = \mathsf{N}_{0}$ of the orbit space, which for $\mathsf{N}_{0} > H^{2}$ has interior points ($|m| < \sqrt{\mathsf{N}_{0}^{2} - H^{4}}$), charge-breaking, and two boundary points, the frames $\pm\eta_{0}$, with no preference: the frame is fixed only through $m$. The conjugation: $c^{*} = \bar{\beta}\cdot v + \bar{\alpha}\cdot w$ exchanges $|\alpha|^{2}$ and $|\beta|^{2}$, so $\mathsf{N} \mapsto \mathsf{N}$, $m \mapsto -m$, $Q \mapsto \bar{Q}$ (checked), and $V_{D}$ with level $\xi$ is carried to $V_{D}$ with level $-\xi$.

\emph{(2).} Substituting (1) into $V_{D}$ gives the displayed form on $L$ (the check verifies it at the off-family point). At fixed $m$ it is strictly increasing in $\mathsf{N} \ge \sqrt{m^{2} + H^{4}}$, so the infimum over the orbit space is attained on the boundary, where $V_{D} = \tfrac18\bigl(g^{2}m^{2} + g'^{2}(m - \xi)^{2}\bigr)$, a strictly convex quadratic in $m$ with minimum at $m^{*} = g'^{2}\xi/(g^{2} + g'^{2})$; since $m = -H^{2}\sinh 2\eta$ is a bijection of the family onto $\R$ (clause (4) of the stored theorem), the vacuum is one orbit, at $\sinh 2\eta^{*} = -m^{*}/H^{2}$, with $\mathsf{N}^{*} = H^{2}\cosh 2\eta^{*} = \sqrt{m^{*2} + H^{4}}$. The check confirms, in three frames with $\xi$ chosen so that the frame is rational, that neighbouring frames have larger $V_{D}$. Nothing in the stored structure fixes $\xi$: the kinetic term, the moment maps and $L$ are all determined, and $\xi$ is the one constant the gauge symmetry allows in the D-term (the Fayet-Iliopoulos constant of the abelian factor, Fact \claimref{synthesis.imported-fayet-iliopoulos-level-only-for-abelian-factors}); so the frame, which is $v/H$ by the premise \claimref{synthesis.rung-unit-second-free-number} and Theorem \claimref{synthesis.rung-unit-is-frame-boost}, is an input of the action. With $\xi = 0$, $m^{*} = 0$ and the vacuum is the balanced frame.

\emph{(3).} \emph{Neutral mode.} Along the family $V_{D}(\eta) = \tfrac18(g^{2}m^{2} + g'^{2}(m - \xi)^{2})$ with $dm/d\eta = -2H^{2}\cosh 2\eta = -2\mathsf{N}$; at the vacuum the first derivative vanishes and $d^{2}V_{D}/d\eta^{2} = \tfrac14(g^{2} + g'^{2})(dm/d\eta)^{2} = (g^{2} + g'^{2})\mathsf{N}^{*2}$. The kinetic coefficient of $(\partial\eta)^{2}$ is $\mathsf{N}(dc_{\lambda}/d\eta) = \mathsf{N}^{*}$ (stored clause (3)), so $m_{h}^{2} = (g^{2} + g'^{2})\mathsf{N}^{*2}/(2\mathsf{N}^{*}) = \tfrac12(g^{2} + g'^{2})\mathsf{N}^{*} = M_{Z}^{2}$. The check does this in the variable $\lambda$: $dm/d\lambda = -H^{4}/\lambda^{3} - \lambda$ (verified as $2G(c, K\,dc/d\lambda)$), $\mathsf{N}(dc/d\lambda) = \mathsf{N}/\lambda^{2}$, and the second derivative along the curve $c_{\lambda^{*}} + \epsilon\,dc/d\lambda + \epsilon^{2}e$ in $L$ obtained by exact interpolation of the degree-eight polynomial $V_{D}(c(\epsilon))$, with the first derivative zero, equals $2\mathsf{N}(d)M_{Z}^{2}$ in all three frames. \emph{Charged modes.} At $c_{\lambda}$ the physical charged plane is $\{\alpha' = a\,e_{5},\ \beta' = (H^{2}/\lambda^{2})\bar{a}\,e_{5} : a \in \C\}$: these $d$ have $Q(c_{\lambda}, d) = 0$ (tangent to $L$), are $G$-orthogonal to $T_{3}c_{\lambda}$, $Yc_{\lambda}$ and $dc_{\lambda}/d\lambda$ (which lie in the $e_{4}$ components) and to $T_{1}c_{\lambda} = (\tfrac{I_{4}H^{2}}{2\lambda}, \tfrac{I_{4}\lambda}{2})e_{5}$, $T_{2}c_{\lambda} = (-\tfrac{H^{2}}{2\lambda}, \tfrac{\lambda}{2})e_{5}$, since $\operatorname{Re}(-I_{4}\tfrac{H^{2}}{2\lambda}a - I_{4}\tfrac{\lambda}{2}\tfrac{H^{2}}{\lambda^{2}}\bar{a}) = 0$ and $\operatorname{Re}(-\tfrac{H^{2}}{2\lambda}a + \tfrac{\lambda}{2}\tfrac{H^{2}}{\lambda^{2}}\bar{a}) = 0$; the check computes this plane as the $G$-orthogonal complement of the gauge orbit in the tangent space of $L$ (three physical directions) with the neutral line removed. For such $d$: $Q(d) = ab = \tfrac{H^{2}}{\lambda^{2}}|a|^{2}$, the second-order correction keeping the curve in $L$ is $e = sc_{\lambda}$ with $s = -Q(d)/2Q(c_{\lambda}) = |a|^{2}/2\lambda^{2}$, real, and $\mathsf{N}(d) = \tfrac12|a|^{2}(1 + H^{4}/\lambda^{4}) = |a|^{2}\mathsf{N}/\lambda^{2}$, $m(d) = -|a|^{2}m/\lambda^{2}$. Along the curve $Q = -H^{2} + O(\epsilon^{3})$, $\mathsf{N} = \mathsf{N}^{*} + \epsilon^{2}n_{2}$ with $n_{2} = \mathsf{N}(d) + 2s\mathsf{N}^{*} = 2|a|^{2}\mathsf{N}^{*}/\lambda^{2}$, and $m = m^{*} + \epsilon^{2}m_{2}$ with $m_{2} = m(d) + 2sm^{*} = 0$ (the first-order terms vanish because $G(c_{\lambda}, d) = m(c_{\lambda}, d) = 0$). Hence $\tfrac12 d^{2}V_{D}/d\epsilon^{2} = \tfrac{g^{2}}{4}\mathsf{N}^{*}n_{2}$ and the mass is $\tfrac{g^{2}}{4}\mathsf{N}^{*}n_{2}/\mathsf{N}(d) = \tfrac12 g^{2}\mathsf{N}^{*} = M_{W}^{2}$, for every $a$. The check verifies it on two basis directions of the plane and on their sum, which determines the quadratic form, in three frames. \emph{Comparison.} The measured values (PDG, from memory) are $m_{h} = 125.20\,\mathrm{GeV}$ and $M_{Z} = 91.19\,\mathrm{GeV}$, so the pure D-term potential is excluded at tree level by a factor $1.37$ in the mass; this is the mechanism of the tree-level bound $m_{h} \le M_{Z}|\cos 2\beta|$ of the supersymmetric two-doublet sector (from memory), with the difference that the constraint $Q(c) = -H^{2}$ removes the second neutral field, so the bound is an equality and no mixing angle appears; what is excluded is the potential, not the locus.

\emph{(4).} Add $W(\mathsf{N})$. On the boundary, $V_{b}(m) = \tfrac18(g^{2}m^{2} + g'^{2}(m - \xi)^{2}) + W(\sqrt{m^{2} + H^{4}})$ with $d\mathsf{N}/dm = m/\mathsf{N}$ and $d^{2}\mathsf{N}/dm^{2} = H^{4}/\mathsf{N}^{3}$, so stationarity is the displayed condition and $V_{b}'' = \tfrac14(g^{2} + g'^{2}) + W''m^{2}/\mathsf{N}^{2} + W'H^{4}/\mathsf{N}^{3}$; with $dm/d\eta = -2\mathsf{N}$ and the kinetic coefficient $\mathsf{N}$, $m_{h}^{2} = 2\mathsf{N} V_{b}'' = M_{Z}^{2} + 2W''m^{*2}/\mathsf{N}^{*} + 2W'H^{4}/\mathsf{N}^{*2}$. For the charged plane the computation of (3) carries over: $\tfrac12 d^{2}V/d\epsilon^{2} = n_{2}(\tfrac{g^{2}}{4}\mathsf{N}^{*} + W'(\mathsf{N}^{*}))$ since $m_{2} = 0$ and $W(\mathsf{N}^{*} + \epsilon^{2}n_{2}) = W + \epsilon^{2}W'n_{2} + O(\epsilon^{4})$, so $m_{\pm}^{2} = 2(\tfrac{g^{2}}{4}\mathsf{N}^{*} + W'(\mathsf{N}^{*})) = M_{W}^{2} + 2W'(\mathsf{N}^{*})$. The outward $\mathsf{N}$-derivative of $V$ on $L$ at fixed $m$ is $\partial_{\mathsf{N}}V = \tfrac{g^{2}}{4}\mathsf{N} + W'(\mathsf{N})$, so $m_{\pm}^{2} = 2\partial_{\mathsf{N}}V$ at the boundary: the family point is a local minimum of $V$ on $L$ iff $V_{b}''(m^{*}) > 0$ (neutral) and $\partial_{\mathsf{N}}V > 0$ (the two charged directions, which are the directions into the interior of the orbit space). This is local stability only: a $W$ that decreases again at larger $\mathsf{N}$ can have a lower charge-breaking minimum in the interior, which the Hessian does not see. The check verifies the three formulas with $W = \kappa(\mathsf{N} - \mathsf{N}_{0})^{2}$ at two frames and two $(\kappa, \mathsf{N}_{0})$, with $\xi$ solved from stationarity, and finds one sample with $m_{\pm}^{2} < 0$ (a charge-breaking instability: $W' < -g^{2}\mathsf{N}^{*}/4$) and one stable. For $v \gg H$, $m^{*2}/\mathsf{N}^{*} = \mathsf{N}^{*} - H^{4}/\mathsf{N}^{*}$, so $m_{h}^{2} = M_{Z}^{2} + 2\mathsf{N}^{*}W''$ up to a relative correction of order $H^{4}/\mathsf{N}^{*2} \approx 10^{-176}$; with $v^{2} = 2\mathsf{N}^{*}$ (the normalisation in which $M_{W}^{2} = g^{2}v^{2}/4$), $m_{h}^{2} = M_{Z}^{2} + v^{2}W''(\mathsf{N}^{*})$, and the measured masses with $v = 246.2\,\mathrm{GeV}$ need $W''(\mathsf{N}^{*}) \approx (125.20^{2} - 91.19^{2})/246.2^{2} \approx 0.121$, that is $\kappa \approx 0.061$ for a quadratic $W$ (numbers in this paragraph only). $W'(\mathsf{N}^{*})$ is not constrained by the neutral mass at this order; it is read off the charged mass, $m_{\pm}^{2} = M_{W}^{2} + 2W'(\mathsf{N}^{*})$. The charged pair is physical (the six real directions of $L$ at the vacuum are three gauge, one neutral and two charged), so the rung field is not a one-doublet Higgs field; in the pure D-term theory the pair is degenerate with the $W$ at about $80\,\mathrm{GeV}$, which direct searches for charged scalars (lower bound near $80\,\mathrm{GeV}$, from memory) place at the edge, so the data constrain $W'(\mathsf{N}^{*}) \ge 0$ as well as $W''(\mathsf{N}^{*})$. The strength $\tfrac12 g^{2}$, $\tfrac12 g'^{2}$ of the D-term relative to the kinetic term is the supersymmetric relation; the gauge symmetry alone allows any positive multiples, and $m_{h} = M_{Z}$ depends on this choice (a multiple $\kappa_{D}$ on both couplings gives $m_{h}^{2} = \kappa_{D}M_{Z}^{2}$ and $m_{\pm}^{2} = \kappa_{D}M_{W}^{2}$ by the same computation).

\emph{Scope.} The action is written on $L$ as a constraint surface, as the stored rung-field theorem does; the kinetic term is the Hermitian one, whose gauge masses the store's two-readings clause records as the electroweak pattern (the bilinear reading gives no scalar potential of this kind, since $Q$ is constant on $L$). Loop corrections, and the second coupling that the $c^{*}$ reading of the up-type masses needs, are not addressed. Transport-convenient: the locus, the family and the frame as $v/H$ are the series'; the D-term, the Fierz identity and the Hessians are standard.
